\section{Role-Based Sovereign Access Control \& Patient Consent}
\label{sec:rbac_consent}

Securing healthcare federations against unauthorized data leakage while facilitating legitimate clinical collaboration requires a defense-in-depth access architecture. This section presents our orthogonal two-tier security enforcement model, defines the consortium role hierarchy, and formalizes the sovereign patient consent engine.

\begin{figure}[t]
\centering
\begin{tikzpicture}[
    scale=0.85, transform shape,
    font=\sffamily\footnotesize,
    node distance=0.7cm and 0.4cm,
    stepBox/.style={
        rectangle, rounded corners=3pt, draw=black, line width=1pt,
        fill=white, text width=7.4cm, align=center, inner sep=4pt
    },
    decisionBox/.style={
        diamond, draw=black, line width=1.1pt, fill=gray!15,
        text width=2.6cm, align=center, inner sep=2pt, aspect=1.8
    },
    allowBox/.style={
        rectangle, rounded corners=3pt, draw=black, line width=1.2pt,
        fill=white, text width=3.4cm, align=center, inner sep=3.5pt, font=\sffamily\scriptsize
    },
    denyBox/.style={
        rectangle, rounded corners=3pt, draw=black, line width=1.2pt, dashed,
        fill=gray!12, text width=3.4cm, align=center, inner sep=3.5pt, font=\sffamily\scriptsize
    },
    patientBox/.style={
        rectangle, rounded corners=4pt, draw=black, line width=1.3pt,
        fill=gray!25, text width=7.4cm, align=center, inner sep=4pt
    },
    flow/.style={
        -{Stealth[length=5pt, width=3.5pt]}, line width=1pt, draw=black
    },
    labelTxt/.style={
        font=\sffamily\scriptsize\bfseries, fill=white, draw=black, inner sep=1.5pt, rounded corners=1pt
    }
]

% Top: Patient Sovereign Policy
\node[patientBox] (patient) {
    \textbf{Patient Sovereign Portal (Asynchronous On-Chain Policy)}\\
    \textit{Interactive Consent Toggles: } $\mathcal{C}(P_k, C_i) \leftarrow \{\text{Granted}, \text{Revoked}\}$\\
    Signed via Patient Private Key $\sigma_{\text{patient}}$ $\bullet$ Committed to $C_{\text{repo}}$
};

% Middle 1: Query Initiated
\node[stepBox, below=0.6cm of patient] (req) {
    \textbf{1. Clinical Traversal Query Initiated}\\
    $\text{BFS-EHR} / \text{DFS-EHR}(P_k, C_i)$ $\bullet$ \textit{Clinician / Specialist Request}
};

% Middle 2: Decision
\node[decisionBox, below=0.7cm of req] (check) {
    \textbf{Consent State?}\\
    $\mathcal{C}(P_k, C_i) \stackrel{?}{=} \text{Granted}$
};

% Bottom: Outcomes side-by-side
\node[allowBox, below left=0.8cm and -0.2cm of check] (grant) {
    \textbf{[+] Access Granted}\\
    1. Retrieve Ciphertext $c$\\
    2. Decrypt with Key $\mathcal{K}$\\
    3. Assert $\mathcal{H}(m) \equiv h$\\
    4. Append to EHR $\mathcal{L}$
};

\node[denyBox, below right=0.8cm and -0.2cm of check] (deny) {
    \textbf{[-] Access Gated / Denied}\\
    $\text{VisibleRecords}(P_k, C_i) \leftarrow \emptyset$\\
    Log Consent Denial on $C_{\text{repo}}$\\
    \textit{Zero PHI Exposure}
};

% Flows
\draw[flow, dashed] (patient) -- node[labelTxt] {Updates Gatekeeper State} (req);
\draw[flow] (req) -- (check);
\draw[flow] (check) -| node[labelTxt, above] {Granted} (grant);
\draw[flow] (check) -| node[labelTxt, above] {Revoked} (deny);

\end{tikzpicture}
\caption{Patient Sovereign Consent Gatekeeper Flowchart. Traversal queries targeting clinical records for patient $P_k$ on chain $C_i$ are dynamically evaluated against the on-chain consent registry $\mathcal{C}(P_k, C_i)$ prior to ciphertext decryption and hash verification.}
\label{fig:consent_gate}
\end{figure}

\subsection{Orthogonal Two-Layer Security Model}
To prevent both presentation-tier bypasses and smart contract privilege escalations, BlockchainForkTree enforces security across two decoupled, complementary architectural layers:
\begin{enumerate}
    \item \textbf{Smart Contract Verification Layer:} At the EVM execution tier, every state-modifying function in \texttt{ConsortiumGovernance.sol} and \texttt{BlockData.sol} verifies the cryptographic identity of the caller (\texttt{msg.sender}). Caller addresses are validated against the organizational registry $\mathcal{O}$ and role authorization mappings stored in the state trie. Any transaction submitted by an unauthorized address immediately triggers an EVM opcode \texttt{REVERT}, reverting all state modifications and consuming minimal execution gas.
    \item \textbf{Presentation-Tier Route Scoping:} At the application programming interface (API) and web interface tier, the Express server (\texttt{server.js}) implements strict cryptographic persona jail-scoping. API routes inspecting inter-organizational message queues and clinical observation stores enforce deterministic query filters parameterized by the authenticated user's bound \texttt{networkId} and role classification. Privileged administrative interfaces (e.g., fork spawning wizards, project creation forms) are completely omitted from the client Document Object Model (DOM) for non-administrative roles, preventing unauthorized interface inspection and reconnaissance.
\end{enumerate}

[Note / Expansion Opportunity: The dual-layer security model can be formalized as a non-interference theorem in programming language theory, proving that lower-privilege personas can never induce state observable by higher-privilege observers.]

\subsection{Formal Persona Permission Matrix}
We categorize consortium stakeholders into six orthogonal personas with strictly delineated operational boundaries:
\begin{itemize}
    \item \texttt{STEERING\_COUNCIL}: Executive governing body authorized to instantiate root federated projects and adjust global protocol parameters.
    \item \texttt{ORG\_ADMIN}: Institutional administrator authorized to submit fork proposals, vote on active network ballots, and register local node daemons.
    \item \texttt{CLINICIAN}: Licensed medical practitioner authorized to author encounters, conditions, and dispatch diagnostic service orders for assigned patients.
    \item \texttt{SPECIALIST}: Laboratory pathologist or pharmacist authorized to inspect assigned diagnostic orders and commit verified clinical reports and medication dispensations.
    \item \texttt{AUDITOR}: Independent regulatory observer possessing read-only privileges to verify on-chain SHA-256 digests and audit topological lineage.
    \item \texttt{PATIENT}: Sovereign health consumer possessing exclusive authority to inspect personal longitudinal health records and dynamically grant or revoke access consent.
\end{itemize}

Table~\ref{tab:rbac} details the complete Role-Based Access Control (RBAC) permission matrix across the platform.

\begin{table*}[t]
\centering
\caption{Role-Based Access Control (RBAC) Persona Permission Matrix}
\label{tab:rbac}
\footnotesize
\setlength{\tabcolsep}{3pt}
\begin{tabularx}{\textwidth}{p{3.8cm}*{6}{>{\centering\arraybackslash}X}}
\toprule
\textbf{Role Classification} & \textbf{Administer Projects} & \textbf{Request/Vote Forks} & \textbf{Author Clinical Records} & \textbf{Fulfill Orders (Lab/Rx)} & \textbf{Verify Audit Hashes} & \textbf{Sovereign Consent} \\
\midrule
\texttt{STEERING\_COUNCIL} & \checkmark & \checkmark & $\times$   & $\times$   & \checkmark & $\times$   \\
\texttt{ORG\_ADMIN}        & $\times$   & \checkmark & $\times$   & $\times$   & \checkmark & $\times$   \\
\texttt{CLINICIAN}        & $\times$   & $\times$   & \checkmark & $\times$   & $\times$   & $\times$   \\
\texttt{SPECIALIST}       & $\times$   & $\times$   & $\times$   & \checkmark & $\times$   & $\times$   \\
\texttt{AUDITOR}          & $\times$   & $\times$   & $\times$   & $\times$   & \checkmark & $\times$   \\
\texttt{PATIENT}          & $\times$   & $\times$   & $\times$   & $\times$   & $\times$   & \checkmark \\
\bottomrule
\end{tabularx}
\end{table*}

[Note / Expansion Opportunity: The RBAC matrix can be augmented with temporal delegation primitives (e.g., temporary 24-hour emergency clinician access) to support urgent trauma care without permanent privilege escalation.]

\subsection{Sovereign Patient Consent Model}
Under HIPAA \S~164.508~\cite{hipaa1996} and GDPR Article 6(1)(a)~\cite{gdpr2016}, cross-organizational processing of Protected Health Information requires demonstrable, freely given, and revocable subject consent. BlockchainForkTree operationalizes this statutory requirement through an on-chain dynamic consent gatekeeper (Fig.~\ref{fig:consent_gate}).

\begin{definition}[Sovereign Consent State Function]
For any registered patient $P_k$ and healthcare institution represented by blockchain $C_i$, the consent state $\mathcal{C}(P_k, C_i) \in \{\text{Granted}, \text{Revoked}\}$ is recorded in an on-chain mapping within \texttt{ConsortiumGovernance.sol}:
\begin{equation}
\mathcal{C}: \mathcal{P}_{\text{patients}} \times V \to \{\text{Granted}, \text{Revoked}\}
\end{equation}
\end{definition}

During longitudinal EHR reconstruction or inter-institutional clinical querying, the resolution function $\text{VisibleRecords}(P_k, C_i)$ dynamically filters the returned observation set:
\begin{equation}
\text{VisibleRecords}(P_k, C_i) = \begin{cases} 
\mathcal{R}(P_k, C_i) & \text{if } \mathcal{C}(P_k, C_i) \equiv \text{Granted} \\
\emptyset & \text{if } \mathcal{C}(P_k, C_i) \equiv \text{Revoked}
\end{cases}
\end{equation}
where $\mathcal{R}(P_k, C_i)$ denotes the set of clinical records for $P_k$ committed on chain $C_i$.

Patients update their consent preferences asynchronously via the web portal by signing an on-chain transaction with their private key, invoking \texttt{setConsent(institutionId, status)}. If a patient revokes consent for Hospital $C_1$, subsequent traversal algorithms executing on behalf of an external provider immediately receive an empty set for $C_1$, while cryptographic audit hashes remain immutably anchored for malpractice verification. This mechanism provides mathematical guarantees of compliance with GDPR Article 7(3) (``The data subject shall have the right to withdraw his or her consent at any time'') without compromising the historical integrity of the underlying ledger.

[Note / Expansion Opportunity: Future work could investigate fine-grained semantic consent policies where patients grant access to specific clinical categories (e.g., cardiology vitals) while revoking sensitive behavioral health encounters.]
